% Generated from paper/full_draft. Edit the Markdown source or migration script.

\section{Method}
\label{sec:04-method:method}

\subsection{Problem Formulation}
\label{sec:04-method:problem-formulation}

In the retrieval-backed implementation evaluated here, let an evidence corpus be a set of chunks $D = \{d_i\}_{i=1}^{N}$ and a query stream $Q = \{q_t\}_{t=1}^{T}$. For each query $q_t$, a retrieval system returns an ordered context $C_t = [d_{t,1}, \ldots, d_{t,k}]$ that will be passed to a downstream generator. The objective is to preserve retrieval quality while controlling route use and the final evidence-input context size.

IntentRoute separates confidence-gated routing from final-context budgeting. For each query, the system estimates how much to rely on global dense retrieval, lexical BM25 recall, and cluster-local retrieval. A separately calibrated policy then budgets the resulting ranked evidence. Retrieval quality and final context tokens remain the primary mechanism-level outcomes; a frozen 300-query generation-and-judge experiment evaluates whether the same trade-off reaches answer-level behavior.

\subsection{Piecewise Relevance-Manifold Assumption}
\label{sec:04-method:piecewise-relevance-manifold-assumption}

The method is motivated by a bounded assumption:

\begin{quote}
In vertical-domain evidence retrieval, query-document relevance often follows a piecewise local structure induced by domain terminology, semantic neighborhoods, document organization, and user intent.
\end{quote}

This does not mean that corpus geometry is sufficient by itself. Instead, IntentRoute uses geometry as one routing signal among several. Dense retrieval remains a global recall floor, BM25 provides lexical anchors, and cluster-local retrieval provides local evidence patches.

\subsection{Multi-Route Retrieval Surface}
\label{sec:04-method:multi-route-retrieval-surface}

IntentRoute uses three retrieval routes.

\textbf{Dense route.} The dense route performs global dense retrieval over the selected corpus using cosine similarity in embedding space. It is the primary quality baseline and the main recall floor.

\textbf{BM25 route.} The BM25 route performs global lexical retrieval. It supports exact-match and terminology-sensitive queries and provides a lexical alternative when dense similarity is insufficient.

\textbf{Cluster-local route.} The cluster-local route performs dense retrieval within LinUCB-selected cluster arms. It allows the system to exploit local structure without searching the entire corpus as aggressively.

The routes are fused in the retrieval layer. The system can run in a full multi-route mode, a gated cost-aware mode, or a final context compaction mode. Full multi-route retrieval improves coverage but does not automatically reduce final context tokens. The stronger token-saving policy described below is calibrated independently after route construction; it does not convert route confidence into a per-query token ratio.

Figure~\ref{fig:system} summarizes this route-control architecture after the route surface is defined. It should be read as a controller diagram rather than a claim that LinUCB replaces dense or BM25 retrieval: dense and BM25 are global recall routes, LinUCB selects cluster-local arms, and the resulting ranking is passed to a separately calibrated final context-budget controller.

\subsection{Routing Arm Construction}
\label{sec:04-method:routing-arm-construction}

Corpus chunk embeddings are clustered with KMeans or MiniBatchKMeans. This is a deliberate experimental choice. LinUCB requires a fixed number of arms, fixed arms improve reproducibility across seeds and scales, and KMeans is fast enough for large-scale LoTTE experiments. The same arm count is used across LoTTE technology/search scales to keep the LinUCB state space comparable, even though larger corpora therefore contain more chunks per arm. We use 32 routing arms as the main reproducible operating point. A sensitivity study over $K \in \{8,16,32,64,128\}$ shows that full multi-route fused quality is stable across this range, while retrieval-stage gated routing is more sensitive because smaller $K$ changes route granularity and larger $K$ spreads feedback across more arms. We therefore treat 32 as an engineering choice, not a manifold-derived optimum.

The paper does not claim that KMeans is the best clustering algorithm for retrieval. HDBSCAN, graph clusters, or learned routing structures may be better in some deployments, but dynamic arm counts complicate the current LinUCB setup. Here, each cluster arm represents a local region of the retrieval surface. The cluster route searches inside selected arms, while dense and BM25 routes can rescue cases where the cluster route misses relevant evidence.

\subsection{LinUCB Route Policy}
\label{sec:04-method:linucb-route-policy}

For each query $q_t$, IntentRoute computes a controller context vector $x_t \in \mathbb{R}^{p}$ by applying a corpus-fitted PCA projection to the query embedding and then L2-normalizing the projected vector. For each arm $a \in \mathcal{A}$, LinUCB maintains a linear value model:

\[
\begin{aligned}
\hat{\theta}_a &= A_a^{-1} b_a, \\
s_t(a) &= \hat{\theta}_a^\top x_t
        + \alpha \sqrt{x_t^\top A_a^{-1} x_t}.
\end{aligned}
\]

The first term estimates the arm value under the current query context. The second term is an exploration bonus. The parameter $\alpha$ controls the exploration-exploitation balance: larger values encourage the policy to explore arms whose value estimate is uncertain.

The policy selects the top candidate arms by score. The selected arms define the cluster-local retrieval path and provide confidence signals for route gating and fallback.

\subsection{Controller Context and Route-Gating Signals}
\label{sec:04-method:controller-context-and-route-gating-signals}

The tested LinUCB context is deliberately narrow: it is the normalized PCA-projected query embedding, not a concatenation of dense scores, BM25 scores, route agreement, feedback summaries, or budget variables. The PCA fit uses corpus embeddings only, so query vectors and fixed KMeans arm centroids share one controller space.

After LinUCB selects arms from this context, the implementation computes policy-derived confidence and a selected-arm centroid mismatch safeguard as separate route-gating signals. With normalized query context $x_t$ and selected arm set $\mathcal{A}_t$, the latter is $1-\max_{a\in\mathcal{A}_t}\cos(x_t,\mu_a)$. It measures mismatch to the chosen centroids; it is not temporal or distribution drift. These signals determine whether the system retains the full dense/BM25/cluster fusion surface or permits a lighter route with a Dense floor. Dense and BM25 rankings are constructed and fused after the route decision; their score concentration, lexical strength, and route overlap are not inputs to the tested LinUCB vector. Neither signal sets the final token ratio, which remains a separately calibrated action.

\subsection{Trust-Weighted Route-State Updates}
\label{sec:04-method:trust-weighted-route-state-updates}

IntentRoute models feedback as a noisy signal rather than a perfect oracle. In the trust-weighted mode, each simulated user feedback event is assigned a trust weight. Higher-trust feedback contributes more to the arm update and local feedback memory. Lower-trust feedback has a weaker effect.

For a selected source arm, let $o_{t,a}$ be the observed simulated-feedback reward and let $w_{t,a}$ be its trust- and propagation-weighted update weight. The tested update is:

\[
\begin{aligned}
A_a &\leftarrow A_a + w_{t,a} x_t x_t^\top, \\
b_a &\leftarrow b_a + w_{t,a} o_{t,a} x_t.
\end{aligned}
\]

The reward is an evidence-quality signal derived from the selected route or the final fused ranking, according to the declared attribution mode. Trust weighting changes the update strength, and neighboring arms can receive a decayed propagated update. Candidate cost and final-context tokens are measured separately; the tested LinUCB reward does not include a $-\lambda c_t$ cost-penalty term.

The current experiments do not claim that real human feedback was collected. They evaluate controlled repeated-interaction adaptation and hard-case recovery. The strongest feedback evidence is visible in policy metrics such as last true reward and selected-cluster hit rate, especially when final retrieval quality is already protected by dense and BM25 fallback.

\subsection{Route-Level Credit Assignment}
\label{sec:04-method:route-level-credit-assignment}

The implementation supports two attribution modes. \nolinkurl{final_fused} assigns reward from the final fused ranking, whereas \nolinkurl{cluster_only} assigns reward from the selected cluster-local route and avoids crediting Dense/BM25 rescue to that arm. The common cross-dataset evidence rows and formal frozen-policy audit use \nolinkurl{final_fused} attribution; \nolinkurl{cluster_only} is retained for dedicated credit-assignment and mechanism diagnostics. Every result family reports its attribution mode in Supplementary Table S22.

This distinction matters because final fused reward measures the system outcome, while cluster-only reward isolates the cluster-route component. Neither mode allows Dense/BM25 rescue to be omitted from the interpretation of final quality.

\subsection{Prequential Adaptation Protocol}
\label{sec:04-method:prequential-adaptation-protocol}

For each query $q_t$, the policy state is frozen before retrieval. The system ranks candidates and evaluates the resulting evidence before converting the observed ground-truth outcome into simulated feedback. The update can therefore change only the route state used by $q_{t+1}$ and later interactions:

\[
\text{rank}(q_t;\theta_t)\rightarrow o_t\rightarrow
\theta_{t+1}=U(\theta_t,x_t,o_t,w_t).
\]

The current query never sees its own label before ranking. Experimental Setup specifies the oracle, noisy, trust-weighted, repeated-interaction, and frozen unseen-query controls used to test this order.

\subsection{Final Context Budgeting and Conservative Confidence Baseline}
\label{sec:04-method:final-context-budgeting-and-conservative-confidence-baseline}

Preliminary token-cost analysis showed that reducing source candidates does not automatically reduce final context tokens. IntentRoute therefore combines confidence-gated routing with an explicit calibrated budget policy, the central quality-cost control interface.

Let $R_t$ be the ranking produced by the confidence-gated route surface. A calibration policy $\pi_\phi$ selects a token ratio $r \in (0,1]$ and a mandatory ranked prefix size $m$. The final context begins with that prefix, then scans the remaining top-10 candidates in rank order and retains a candidate only when it fits the residual budget:

\[
\mathrm{Tokens}(C_t) \le r\,\mathrm{Tokens}(R_t[:10]).
\]

This is an order-preserving budgeted subset with a mandatory prefix. It may skip an oversized later chunk and admit a smaller lower-ranked chunk, so it is not necessarily a contiguous longest ranked prefix. The policy parameters $\phi$ are selected on calibration queries subject to a retrieval-quality eligibility gate and then frozen before held-out test evaluation. Geometry and feedback affect route construction, while the stronger token saving arises when the separately calibrated length budget acts on the routed ranking. The implementation does not learn a direct per-query mapping from confidence to token ratio.

The conservative \nolinkurl{confidence_topk} policy works as follows:

\begin{enumerate}
\item If route confidence is low or selected-arm centroid mismatch is high, keep dense fallback and the normal top-10 context.
\item If LinUCB confidence is high, reduce final context size from the default top-10 to $k=8$.
\item If confidence is mid-level, keep $k=10$ as a safety tier rather than calling it true compression.
\item Report the actual retrieved context token count, not only candidate counts.
\end{enumerate}

For the conservative policy, the effective compaction rate is therefore the high-confidence compression rate, not the broader rate of queries that enter a non-fallback route. Hybrid-lite can reduce dense influence in fusion while retaining dense candidates as a safety net; it should not be described as reducing dense computation unless the global dense route is actually skipped.

The main token-quality result uses frozen calibration/test budget policies. They select a global ratio and minimum-prefix size on calibration queries, then apply both unchanged to held-out test queries. The conservative confidence-based policy remains a stable baseline: it reduces final context tokens by about 4.7-5.3\% across LoTTE technology/search at 100k, 200k, 400k, and 638k while preserving dense-level query hit.

For the conservative \nolinkurl{confidence_topk} baseline only, centroid mismatch triggers fallback in roughly 1--2\% of interactions under the configured threshold, so context-size decisions are primarily confidence-conditioned. The fixed-pool factorial audit does not show that this confidence predicts compression safety better than matched controls; the baseline is an empirical operating point rather than a validated causal mechanism.

\subsection{Feedback-Triggered Recovery}
\label{sec:04-method:feedback-triggered-recovery}

IntentRoute can also use feedback as a recovery signal after a compressed context fails. In this mode, feedback does not insert the missing ground-truth chunk into the current answer. Instead, it updates arm-level route value and marks the associated local region as risky for future routing or optional post-feedback retry.

The safe recovery behavior is conservative:

\begin{enumerate}
\item If feedback indicates that a compressed context missed evidence, identify the evidence-bearing arm or risky query-arm relation.
\item Prefer a less aggressive final-context budget or full-context fallback for that local region.
\item Avoid unconditional global arm boosting, because a feedback-positive arm for one query can transfer poorly to unrelated queries.
\end{enumerate}

This recovery mechanism is deployment-facing. It represents how a system can repair tail failures after user or evaluator feedback, not how the first-pass ranking is produced before feedback is observed.

\subsection{Algorithm Sketch}
\label{sec:04-method:algorithm-sketch}

Input: query $q_t$, corpus $D$, route artifacts, and the current LinUCB state.

\begin{enumerate}
\item Embed $q_t$ and form the normalized PCA controller context $x_t$.
\item Score cluster arms with LinUCB using $s_t(a)$.
\item Retrieve candidates from the global dense route, the global BM25 route, and dense search within selected cluster arms.
\item Fuse route rankings.
\item Apply the independently calibrated final-context budget, or the explicitly labeled conservative \nolinkurl{confidence_topk} baseline.
\item Evaluate retrieval quality for $q_t$.
\item Only after evaluation, convert the ground-truth label into simulated feedback and update LinUCB for later queries.
\item If the interaction is a post-feedback retry or a later query in a risky local region, use the updated arm state to choose a safer route or fallback path; final-context budgeting remains independently calibrated.
\end{enumerate}

Output: final retrieved context $C_t$ and updated policy state.

\subsection{Reproducibility Parameters}
\label{sec:04-method:reproducibility-parameters}

The common cross-dataset evidence protocol uses 32 fixed KMeans/MiniBatchKMeans arms, three candidate arms per query, a 64-dimensional context projection, $\alpha=1.0$, and eight prequential epochs. Earlier mechanism diagnostics may use other epoch counts and attribution modes; Supplementary Table S22 records those result-family differences. The calibration grid covers $r \in \{0.85,0.88,0.90,0.92,0.95,0.98\}$ and $m \in \{4,\ldots,8\}$. Supplementary Section S12 reports the complete route depths, safety floors, fusion weights, confidence thresholds, decay, and trust parameters; scale-specific cache paths and dataset sizes remain in the tracked experiment artifacts.

The notation \nolinkurl{token_budget_r0.85_m4} means that each query keeps a safe prefix of at least four chunks and then admits additional chunks only while the final context remains within 85\% of the original dense top-10 token budget. The policy is chosen on calibration queries and then frozen before held-out test evaluation.
